% uid: ""
% school: "lycee"
% level: "3eme"
% section: ["Economie et Gestion"]
% chapter: "Statistiques"
% title: "Nombre d'enfants par famille"
% difficulty: "Facile"
% points: null
% tags: ["statistiques", "tableau-d-effectifs", "moyenne", "ecart-type", "quartiles"]
% provenance: "own"
% dressing: "tableau d'effectifs"
% rung: 2
% batch: "2026-11-20-stat-3eco-A"

% ---STATEMENT---
On a interrogé $50$ familles d'un quartier sur leur nombre d'enfants $x$ :

\begin{center}
\begin{tabular}{|l|c|c|c|c|c|c|}
\hline
Nombre d'enfants $x_i$ & $0$ & $1$ & $2$ & $3$ & $4$ & $5$ \\ \hline
Effectif $n_i$ & $6$ & $14$ & $15$ & $9$ & $4$ & $2$ \\ \hline
\end{tabular}
\end{center}

\begin{enumerate}
  \item Calculer le nombre moyen d'enfants par famille.
  \item Calculer la variance puis l'écart type (arrondi à $10^{-2}$).
  \item Déterminer la médiane et les quartiles $Q_1$, $Q_3$.
  \item Quel pourcentage de familles a au plus $2$ enfants ?
\end{enumerate}

% ---SOLUTION---
On commence par compléter le tableau de calcul.
\begin{center}
\small
\begin{tabular}{|l|c|c|c|c|c|c|c|}
\hline
$x_i$ & $0$ & $1$ & $2$ & $3$ & $4$ & $5$ & Total \\ \hline
$n_i$ & $6$ & $14$ & $15$ & $9$ & $4$ & $2$ & $50$ \\ \hline
$n_ix_i$ & $0$ & $14$ & $30$ & $27$ & $16$ & $10$ & $97$ \\ \hline
$n_ix_i^2$ & $0$ & $14$ & $60$ & $81$ & $64$ & $50$ & $269$ \\ \hline
Effectif cumulé croissant & $6$ & $20$ & $35$ & $44$ & $48$ & $50$ & \\ \hline
\end{tabular}
\end{center}
\begin{enumerate}
  \item \textbf{Rappel.}\par
  \begin{prop}[Moyenne d'une série à effectifs]
  $\displaystyle\overline{x}=\frac{1}{N}\sum n_ix_i$.
  \end{prop}
  $$\overline{x}=\frac{97}{50}=\ib{$1{,}94$}\text{ enfant par famille.}$$
  \item $\displaystyle\frac1N\sum n_ix_i^2=\frac{269}{50}=5{,}38$, donc
  $$V=5{,}38-1{,}94^2=5{,}38-3{,}7636=1{,}6164\qquad\sigma=\sqrt{1{,}6164}\approx\ib{$1{,}27$}.$$
  \item $N=50$.

  \emph{Médiane} : $N$ pair, termes de rang $25$ et $26$. Les effectifs cumulés $20<25$ et $35\ge26$ montrent que ces deux termes valent $2$ : $M_e=2$.

  \emph{$Q_1$} : $\dfrac{50}4=12{,}5\to$ rang $13$. Effectifs cumulés croissants : $6<13$ pour $x=0$, puis $20\ge13$ pour $x=1$ : $Q_1=1$.

  \emph{$Q_3$} : $\dfrac{3\times50}4=37{,}5\to$ rang $38$. Effectifs cumulés croissants : $35<38$ pour $x=2$, puis $44\ge38$ pour $x=3$ : $Q_3=3$.
  \item Au plus $2$ enfants signifie $x\le2$ : c'est l'effectif cumulé en $2$, soit $35$ familles.
  $$\frac{35}{50}=0{,}70=\ib{$70\,\%$}.$$
\end{enumerate}

\textbf{Idée :} un tableau d'effectifs est une série où chaque valeur est répétée : on pondère chaque valeur par son effectif.
